\section{现金流记录与用途相互独立}

本案当前没有获批市场组合，因此没有由本框架指定的账户现金动作。本章只定义已结算分配和外部现金流的记录、归因与边界，绝不构成买入、留存、取出或再配置指令。现金到账不因付款方、除息日、历史分配率或回报目标而触发自动动作；现金用途由投资者独立决定。若日后在独立授权下考虑新的配置，必须作为带日期的新组合审议并通过全部有效资格闸门。

\begin{exhibitbox}[已结算现金：记录、归因与用途边界]
\small
\begin{tabularx}{\exhibitboxwidth}{L{2.2cm}P P}
\toprule
\textbf{步骤} & \textbf{研究动作} & \textbf{控制点} \\
\midrule
1. 入账与对账 & 记录到账日、来源、税前金额、税额、费用、净额、结算凭证及返还性成分标记。 & 公告、应收款和卖出所得不能提前或错误归类为分配；返还性成分不得被称为收益。 \\
2. 总回报归因 & 以统一台账将已结算分配、适用税费与成本纳入TWRR核算一次。 & 留存、取出或后续再配置均不增加第二笔回报，也不改变既有归因。 \\
3. 用途自主 & 投资者可按自身账户、税务、流动性与授权条件决定保留、取出或不作处理。 & 本框架不设默认用途，不把到账现金视为自动配置权限。 \\
4. 另行审议 & 如投资者独立考虑新的配置，才可提交T3候选的边际价值比较。 & $q_{portfolio}$、风险、估值、税费、容量、集中度或相关性任一未通过即不得形成市场配置。 \\
\bottomrule
\end{tabularx}
\end{exhibitbox}

\section{总回报验收、归因与失败状态}

在具有完整实际账户记录、冻结组合及冻结比较基准后，才可做策略业绩验收。总回报采用TWRR；投资者体验另报MWRR/XIRR。正式回放和后续业绩报告均应展示年度与滚动五年TWRR、成本、税费、回撤、波动、下行偏差、换手、未成交、容量、相关性和集中度，并把分配和公司行动纳入同一总回报口径。首次完整五年观察期前，任何“达到／未达到5\%”的判断都应标注为\textbf{尚无完整观察期}；候选基金或单项目历史不能替代策略结果。

\begin{exhibitbox}[实施后才会使用的总回报状态]
\small
\begin{tabularx}{\exhibitboxwidth}{L{3.0cm}P P}
\toprule
\textbf{结果} & \textbf{报告标签} & \textbf{治理结果} \\
\midrule
完整数据下滚动五年TWRR达到或高于5\% & 条件性研究目标达到 & 仍须重做风险、基准、容量和资格审查；不自动扩张风险。 \\
完整数据下滚动五年TWRR低于5\% & 条件性研究目标未达到 & 对账市场／NAV、公司行动、成本、税费、运营与风险暴露；不得通过提高收益风险来“修复”结果。 \\
观察期不足或关键输入缺失 & 尚无完整观察期／不可判定 & 补齐点时序列、成本税费、公司行动和可复现规则；不以单基金历史代替策略回放。 \\
\bottomrule
\end{tabularx}
\end{exhibitbox}

若总回报证据、风险预算、集中度、适当性或实施资格任一项失效，冻结任何增加市场风险的提案，直至相关证据重新通过。这个治理动作不承诺通过交易“修复”结果，也不对现金分配设置独立的年度成功或失败标准。
